\chapter{总结展望}
\label{chap:conclusion}

\section{本文工作总结}

本文面向智能体大语言模型在线上持续部署场景下"按能力域分批、顺序训练"的现实约束，设计并实现了一套完整的持续强化学习训练与评测系统。

在\textbf{方法层面}，本文提出了三个核心组件：（1）一个面向持续学习的复合训练目标，在 GRPO 策略梯度之外叠加逆向 KL 锚定、经验回放损失与固定熵正则，以零框架改动的方式注入；（2）一个按能力分区的九桶经验回放缓冲区，桶固定容量由平方根加权分配、桶内淘汰禁跨桶挤占、桶间按能力空间距离采样回放；（3）一套奖励由外部冻结裁判统一打分的多维乘性聚合，安全作为二元门控提高评价一致性。

在\textbf{系统层面}，本文在不 fork 底层强化学习框架的前提下，通过外部注入接口与钩子将持续学习逻辑挂入训练链路，推理与采样采用完全异步的策略分离部署，并额外落盘框架未提供的离散度指标以监控训练健康度。

在\textbf{评测层面}，本文提出一套按桶分组的抗遗忘评测协议——训完全部能力域后一次评测、权重后置——在省算力的同时量化"越早训练的桶遗忘越重"这一持续学习核心指标。

\section{不足与展望}

本文工作存在以下可改进方向：

\begin{enumerate}
  \item \textbf{桶内淘汰算法有待优化}。当前采用 FIFO（先进先出），是一种最简单的淘汰策略。更具信息量的方案（如按优先级淘汰）在更复杂的训练动态下可能更有效，但需配套更精细的优先级函数以避免系统性地挤出某类样本。

  \item \textbf{桶间距离度量待实证}。当前桶间距离采样基于能力空间的固定嵌入坐标，该坐标反映的是基准任务的类别邻近关系，而非策略遗忘的动力学邻近关系（即"训桶 A 后桶 B 更易被遗忘"的经验因果）。通过实际训练的遗忘模式反推距离坐标，可能比固定坐标更准。

  \item \textbf{U 形块加权暂未开启}。本文将 U 形加权（$\gamma=\delta=1$→退化为均权）设为默认关闭，作为独立消融项保留。其边际贡献是否显著——能否把信用集中到早期分支决策与最终答案段——有待实验验证。

  \item \textbf{基座规模待扩展}。当前验证基于 Qwen3.5-9B，方法设计面向更大规模（如 27B）。不同规模下遗忘模式是否一致、缓冲区容量是否需要按比例缩放，是后续迁移需回答的问题。

  \item \textbf{多模态能力域}。本文聚焦纯文本能力域，未纳入多模态任务。后者因其模态差异与数据稀疏性，可能需要分立的桶策略或不同的回放频率，是自然延伸方向。
\end{enumerate}
